%! Principles of Experimental Design — Unit 1, Lesson 1.6 — Warm-Up (Answer Key)
% PILOT SHAPE (2026-09-06): modeled on AP Statistics lesson 1.4 minus its AP Practice page.
% ONE PAGE, blank and key. Three items at 12pt, \workrowsep 50pt, item spacing 22pt, so each
% item gets real handwriting room and still fits. Do not add a fourth item — the page is full.
% GRADUAL RELEASE: the warm-up rehearses the arithmetic every row of the notes runs on, on the
% very experiment the notes open with. Items 1 and 2 are the two group rates (45/60 = 75%,
% 33/60 = 55%) — equal denominators, because a generator built the groups. Item 3(a) is the
% 20-point gap; item 3(b) reaches back to Lesson 1.5's WATCHED result on the same Study Center
% (65% vs 80%, a 15-point gap the other way), so both gaps are on the board before the hook
% asks which study to believe — and problem 14 (the crux) can ask whether gap size is the reason.
\documentclass[12pt]{article}
\usepackage{saar-article}
\usepackage{saar-key}   % pulls in -boxes; do NOT also load -boxes
\newcommand{\TallMath}[1]{$\displaystyle #1\rule[-1.4em]{0pt}{3.2em}$}
% Word bank strip for every fill-in cluster (user request 2026-09-06). Defined per document;
% the key inherits it and prints the same strip, so heights match.
\newcommand{\wordbank}[1]{\par\vspace{2pt}\noindent\fcolorbox{goldacc}{goldbg}{\parbox{\dimexpr\linewidth-2\fboxsep-2\fboxrule\relax}{\small\textbf{Word bank:} #1}}\par\vspace{4pt}}
% Handwriting room inside every \begin{work} block. Moves the blank and the
% key together, so the two can never drift apart.
\setlength{\workrowsep}{50pt}

\begin{document}

\pageheader{Unit 1, Lesson 1.6}{Warm-Up --- Answer Key}
% NAMESTRIP: no \namedateperiod here — mirrors the blank. NO teachernote here —
% teacher prose lives in the lesson plan.

\vspace{0.15in}

\begin{notesbox}{Spiral Review --- Riverbend High School, This Time They Assign}
\normalsize
Riverbend High School asked for volunteers to test its after-school Study Center.
\textbf{$120$ students volunteered.} A random number generator split them into
\textbf{two groups of $60$}: one group attended twice a week, the other did not
attend at all. At the end of the quarter the counselor counted how many in each
group passed every class.

\begin{enumerate}[leftmargin=1.8em, itemsep=22pt, topsep=10pt]

  \item $45$ of the $60$ students who attended passed every class. What
        \textbf{percent} of that group is it?

        \begin{work}
          \dfrac{45}{60} &= 0.75 \\
                         &= 75\%
        \end{work}

  \item $33$ of the $60$ students who did not attend passed every class. What
        \textbf{percent} of that group is it?

        \begin{work}
          \dfrac{33}{60} &= 0.55 \\
                         &= 55\%
        \end{work}

  \item Compare the two studies of the same Study Center.

        \textbf{(a)} This quarter, how many \textbf{percentage points} higher is
        the group that attended?

        \begin{work}
          75 - 55 &= 20 \text{ percentage points}
        \end{work}

        \textbf{(b)} Last quarter (Lesson 1.5) the counselor only \emph{watched}:
        $65\%$ of the students who used the Study Center passed every class,
        against $80\%$ of the students who never went. Which group was higher
        then, and by how many percentage points? \ans{the non-users, by $80 - 65 = 15$ points}

\end{enumerate}
\end{notesbox}

\end{document}
